\subsubsection{Market Impact Models}

\begin{taxonomy} \hlt{Model Family Segmentation}
\begin{enumerate}[label=\roman*.]
\setlength{\itemsep}{0pt}
\item \hlt{Almgren-Chriss}: optimise a schedule by balancing expected impact cost against execution risk;
\item \hlt{linear impact}: use tractable instantaneous or temporary cost proportional to trading rate or signed size;
\item \hlt{propagator}: model each signed trade as a decaying impulse in future prices;
\item \hlt{empirical square-root}: summarise metaorder cost as concave in order size relative to market volume.
\end{enumerate}
\end{taxonomy}

\begin{principle} \hlt{Permanent Impact or Price Anticipation}\\
Long-horizon impact mixes metaorder pressure and information revealed by trading. Remove the information effect before calling the residual permanent impact.
\end{principle}

\begin{taxonomy} \hlt{Permanent-Impact View Segmentation}
\begin{enumerate}[label=\roman*.]
\setlength{\itemsep}{0pt}
\item \hlt{mechanical view}: prices move because buying and selling pressure are imbalanced;
\item \hlt{informational view}: prices move because the trade reveals new expectations about value;
\item \hlt{hybrid view}: executions contain both pressure and information, so identification matters more than labels.
\end{enumerate}
\end{taxonomy}

\begin{process} \hlt{Debiasing Long-Horizon Impact}
\begin{enumerate}[label=\roman*.]
\setlength{\itemsep}{0pt}
\item classify trades by whether they are likely information-driven or liquidity-driven;
\item remove market and sector moves with a benchmark or factor model;
\item adjust for autocorrelated metaorder flow after the execution day;
\item measure the remaining post-execution drift over several days.
\end{enumerate}
\end{process}

\begin{model} \hlt{CAPM-Like Post-Execution Decomposition}\\
For event day $D$, a stock's post-execution move can be split into systematic and idiosyncratic parts:
\begin{equation*}
\log(P_d)-\log(P_{d-1})=
\beta\left[\log(I_d)-\log(I_{d-1})\right]+\Delta W_d.
\end{equation*}
Here $P_d$ is the stock close, $I_d$ the reference index, $\beta$ the index exposure, and $\Delta W_d$ the residual stock-specific move. Cumulative $\Delta W_d$ isolates the profile after removing index co-movement.
\end{model}

\begin{caveat} \hlt{Autocorrelated Flow Bias}\\
Metaorder flow can remain signed and autocorrelated for days. If later same-direction executions are ignored, their impact may be falsely attributed to the original order.
\end{caveat}
